\documentclass[10pt]{article}
\usepackage[T1]{fontenc}
\usepackage[utf8]{inputenc}
\usepackage[french]{babel}
\usepackage{amsmath,amssymb,mathrsfs,mathpazo}
\usepackage{geometry}
\geometry{a4paper,margin=2cm}
\usepackage{graphicx,array,tabularx,tabularray,colortbl,multirow,multicol,booktabs,enumitem,pifont,ulem,fancybox,diagbox,cancel}
\usepackage{pstricks,pstricks-add,pst-plot,pst-tree,pst-node,pst-text,pst-func,pst-math,pst-bspline,pst-3dplot,pst-eucl,pst-solides3d}
\let\pscalculate\undefined
\usepackage{pst-bezier}
\usepackage{tikz,pgfplots,tkz-tab}
\usetikzlibrary{arrows,arrows.meta,patterns,calc,positioning,trees,decorations.markings,shapes.geometric,shapes.misc,shapes.arrows,intersections,decorations.pathmorphing,decorations.pathreplacing,backgrounds,shadows,babel}
\pgfplotsset{compat=1.18}
\usepgfplotslibrary{fillbetween}
\usepackage[np]{numprint}
\usepackage{listings,makecell,pas-tableur,twoopt}
\let\barre\undefined
\usepackage[pstricks]{tablvar}
\definecolor{darkgreen}{RGB}{0,100,0}
\providecommand{\red}{\color{red}}
\providecommand{\blue}{\color{blue}}
\providecommand{\green}{\color{green}}
\providecommand{\black}{\color{black}}
\providecommand{\white}{\color{white}}
\providecommand{\gray}{\color{gray}}
\providecommand{\cyan}{\color{cyan}}
\providecommand{\magenta}{\color{magenta}}
\providecommand{\yellow}{\color{yellow}}
\providecommand{\violet}{\color{violet}}
\providecommand{\orange}{\color{orange}}
\providecommand{\euro}{\texteuro}
\providecommand{\EUR}{\texteuro}
\providecommand{\ieme}{\textsuperscript{e}}
\providecommand{\degres}{\ensuremath{^\circ}}
\providecommand{\degree}{\ensuremath{^\circ}}
\providecommand{\vv}[1]{\overrightarrow{#1}}
\providecommand{\Coord}[1]{#1}
\providecommand{\decosix}{}
\providecommand{\starredbullet}{\ensuremath{\bullet}}
\pagestyle{empty}
\setlength{\parindent}{0pt}
\usepackage{adjustbox}\DeclareUnicodeCharacter{2236}{:}\DeclareUnicodeCharacter{2190}{\ensuremath{\leftarrow}}
\providecommand{\R}{\mathbb{R}}
\providecommand{\N}{\mathbb{N}}
\providecommand{\Z}{\mathbb{Z}}
\providecommand{\Q}{\mathbb{Q}}
\providecommand{\D}{\mathbb{D}}
\providecommand{\C}{\mathbb{C}}
\providecommand{\e}{{\mathrm{e}}}
\providecommand{\vect}[1]{\overrightarrow{#1}}
\providecommand{\vectt}[1]{\overrightarrow{\mathrm{#1}}}
\providecommand{\vvt}[1]{\overrightarrow{\mathrm{#1}}}
\providecommand{\np}[2][]{{#2~#1}}
\providecommand{\vv}[1]{{\overrightarrow{#1}}}
\renewcommand{\d}{{\mathrm{d}}}
\providecommand{\degres}{{^\circ}}
\providecommand{\degre}{{^\circ}}
\providecommand{\blue}{\color{blue}}
\providecommand{\red}{\color{red}}
\providecommand{\enskip}{\quad}
\providecommand{\llbracket}{\lbrack\!\lbrack}
\providecommand{\rrbracket}{\rbrack\!\rbrack}
\providecommand{\euro}{{\text{€}}}
\providecommand{\no}{{\text{n°}}}
\providecommand{\interval}[2]{{[#1;#2]}}
\providecommand{\intervFO}[2]{{[#1;#2[}}
\providecommand{\dotfill}{\ldots}

\providecommand{\ds}{\displaystyle}
\providecommand{\textding}[1]{#1}
\providecommand{\Oij}{\(\mathrm{O};\vec{\imath},\vec{\jmath}\)}
\providecommand{\Oijk}{\(\mathrm{O};\vec{\imath},\vec{\jmath},\vec{k}\)}
\providecommand{\Ouv}{\(\mathrm{O};\vec{u},\vec{v}\)}
\providecommand{\intervFF}[2]{\ensuremath{[#1;#2]}}
\renewcommand{\euro}{\ifmmode\text{\texteuro}\else\texteuro\fi}
\DeclareUnicodeCharacter{00B0}{\ifmmode^{\circ}\else\textdegree\fi}
\DeclareUnicodeCharacter{20AC}{\ifmmode\text{\texteuro}\else\texteuro\fi}
% Énoncés : Éducation nationale. Transcription et corrigés : APMEP, Denis Vergès et contributeurs. Découpage et classement : CoopMaths, Sébastien Lozano, Éric Elter et contributeurs. Adaptation de la présentation : Maths Explorer.
% CC BY-SA : https://coopmaths.fr/www/about/
% https://coopmaths.fr/alea/static/bac/2023/tex/bac_2023_03_sujet1_ameriquenord_2.tex
\begin{document}
\noindent\textit{APMEP / CoopMaths. CC BY-SA. Présentation adaptée par Maths Explorer.}\par\medskip
\medskip

On considère la fonction $f$ définie sur $\R$ par :

\[f(x)=\mathrm{e}^{3 x}-(2 x+1) \mathrm{e}^{x}\]

Le but de cet exercice est d'étudier la fonction $f$ sur $\mathbb{R}$.

\medskip

\textbf{Partie A - Étude d'une fonction auxiliaire}

\medskip

On définit la fonction $g$ sur $\mathbb{R}$ par :

\[g(x)=3 \mathrm{e}^{2 x}- 2 x - 3\]

\begin{enumerate}
\item 
	\begin{enumerate}
		\item Déterminer la limite de la fonction $g$ en $-\infty$.
		\item Déterminer la limite de la fonction $g$ en $+\infty$.
	\end{enumerate}
\item 
	\begin{enumerate}
		\item On admet que la fonction $g$ est dérivable sur $\mathbb{R}$, et on note $g'$ sa fonction dérivée. Démontrer que pour tout nombre réel $x$, on a $g'(x)=6 \mathrm{e}^{2 x}-2$.
		\item Étudier le signe de la fonction dérivée $g'$ sur $\mathbb{R}$.
		\item En déduire le tableau de variations de la fonction $g$ sur $\mathbb{R}$. Vérifier que la fonction $g$ admet un minimum égal à $\ln (3) - 2$.
	\end{enumerate}
\item 
	\begin{enumerate}
		\item Montrer que $x=0$ est solution de l'équation $g(x) = 0$.
		\item Montrer que l'équation $g(x)=0$ admet une deuxième solution, non nulle, notée $\alpha$, dont on donnera un encadrement d'amplitude $10^{-1}$.
	\end{enumerate}
\item Déduire des questions précédentes le signe de la fonction $g$ sur $\mathbb{R}$.
	\end{enumerate}

\medskip

\textbf{Partie B - Étude de la fonction \boldmath$f$\unboldmath}

\medskip

\begin{enumerate}
\item La fonction $f$ est dérivable sur $\mathbb{R}$, et on note $f'$ sa fonction dérivée.

Démontrer que pour tout nombre réel $x$, on a $f'(x)=\mathrm{e}^{x} g(x)$, où $g$ est la fonction définie dans la \textbf{partie A}.
\item En déduire alors le signe de la fonction dérivée $f'$ puis les variations de la fonction $f$ sur $\mathbb{R}$.
\item Pourquoi la fonction $f$ n'est-elle pas convexe sur $\mathbb{R}$ ? Expliquer.
\end{enumerate}

\bigskip

\end{document}
